\documentclass{report}
\usepackage{hyperref}
% WARNING: THIS SHOULD BE MODIFIED DEPENDING ON THE LETTER/A4 SIZE
\oddsidemargin 0cm
\evensidemargin 0cm
\marginparsep 0cm
\marginparwidth 0cm
\parindent 0cm
\setlength{\textwidth}{\paperwidth}
\addtolength{\textwidth}{-2in}


% Conditional define to determine if pdf output is used
\newif\ifpdf
\ifx\pdfoutput\undefined
\pdffalse
\else
\pdfoutput=1
\pdftrue
\fi

\ifpdf
  \usepackage[pdftex]{graphicx}
\else
  \usepackage[dvips]{graphicx}
\fi

% Write Document information for pdflatex/pdftex
\ifpdf
\pdfinfo{
 /Author     (Pasdoc)
 /Title      ()
}
\fi


% definitons for warning and note tag
\usepackage[most]{tcolorbox}
\newtcolorbox{tcbwarning}{
 breakable,
 enhanced jigsaw,
 top=0pt,
 bottom=0pt,
 titlerule=0pt,
 bottomtitle=0pt,
 rightrule=0pt,
 toprule=0pt,
 bottomrule=0pt,
 colback=white,
 arc=0pt,
 outer arc=0pt,
 title style={white},
 fonttitle=\color{black}\bfseries,
 left=8pt,
 colframe=red,
 title={Warning:},
}
\newtcolorbox{tcbnote}{
 breakable,
 enhanced jigsaw,
 top=0pt,
 bottom=0pt,
 titlerule=0pt,
 bottomtitle=0pt,
 rightrule=0pt,
 toprule=0pt,
 bottomrule=0pt,
 colback=white,
 arc=0pt,
 outer arc=0pt,
 title style={white},
 fonttitle=\color{black}\bfseries,
 left=8pt,
 colframe=yellow,
 title={Note:},
}

\begin{document}
\label{toc}\tableofcontents
\newpage
% special variable used for calculating some widths.
\newlength{\tmplength}
\chapter{Unit ok{\_}unicode{\_}identifiers{\_}utf8}
\label{ok_unicode_identifiers_utf8}
\index{ok{\_}unicode{\_}identifiers{\_}utf8}
\section{Description}
Delphi allows any Unicode character in an identifier, see \href{https://stackoverflow.com/questions/62143736/what-unicode-characters-can-be-used-in-delphi-source-code}{https://stackoverflow.com/questions/62143736/what-unicode-characters-can-be-used-in-delphi-source-code} \href{https://docwiki.embarcadero.com/RADStudio/Athens/en/Identifiers}{https://docwiki.embarcadero.com/RADStudio/Athens/en/Identifiers}

Here we save file with UTF{-}8 encoding and some German characters.
\section{Overview}
\begin{description}
\item[\texttt{\begin{ttfamily}TMyClass\end{ttfamily} Class}]
\item[\texttt{\begin{ttfamily}TfrmArbeitsÜbersicht\end{ttfamily} Class}]
\end{description}
\section{Classes, Interfaces, Objects and Records}
\ifpdf
\subsection*{\large{\textbf{TMyClass Class}}\normalsize\hspace{1ex}\hrulefill}
\else
\subsection*{TMyClass Class}
\fi
\label{ok_unicode_identifiers_utf8.TMyClass}
\index{TMyClass}
\subsubsection*{\large{\textbf{Hierarchy}}\normalsize\hspace{1ex}\hfill}
TMyClass {$>$} TObject
%%%%Description
\subsubsection*{\large{\textbf{Fields}}\normalsize\hspace{1ex}\hfill}
\begin{list}{}{
\settowidth{\tmplength}{\textbf{pmEntLöchen}}
\setlength{\itemindent}{0cm}
\setlength{\listparindent}{0cm}
\setlength{\leftmargin}{\evensidemargin}
\addtolength{\leftmargin}{\tmplength}
\settowidth{\labelsep}{X}
\addtolength{\leftmargin}{\labelsep}
\setlength{\labelwidth}{\tmplength}
}
\label{ok_unicode_identifiers_utf8.TMyClass-pmEntLöchen}
\index{pmEntLöchen}
\item[\textbf{pmEntLöchen}\hfill]
\ifpdf
\begin{flushleft}
\fi
\begin{ttfamily}
public pmEntLöchen: TPopupMenu;\end{ttfamily}

\ifpdf
\end{flushleft}
\fi


\par  \label{ok_unicode_identifiers_utf8.TMyClass-lbledtHöhe}
\index{lbledtHöhe}
\item[\textbf{lbledtHöhe}\hfill]
\ifpdf
\begin{flushleft}
\fi
\begin{ttfamily}
public lbledtHöhe: TLabeledEdit;\end{ttfamily}

\ifpdf
\end{flushleft}
\fi


\par  \label{ok_unicode_identifiers_utf8.TMyClass-btnAufträge}
\index{btnAufträge}
\item[\textbf{btnAufträge}\hfill]
\ifpdf
\begin{flushleft}
\fi
\begin{ttfamily}
public btnAufträge: TButton;\end{ttfamily}

\ifpdf
\end{flushleft}
\fi


\par  \end{list}
\ifpdf
\subsection*{\large{\textbf{TfrmArbeitsÜbersicht Class}}\normalsize\hspace{1ex}\hrulefill}
\else
\subsection*{TfrmArbeitsÜbersicht Class}
\fi
\label{ok_unicode_identifiers_utf8.TfrmArbeitsÜbersicht}
\index{TfrmArbeitsÜbersicht}
\subsubsection*{\large{\textbf{Hierarchy}}\normalsize\hspace{1ex}\hfill}
TfrmArbeitsÜbersicht {$>$} TObject
%%%%Description
\subsubsection*{\large{\textbf{Fields}}\normalsize\hspace{1ex}\hfill}
\begin{list}{}{
\settowidth{\tmplength}{\textbf{pmEntLöchen}}
\setlength{\itemindent}{0cm}
\setlength{\listparindent}{0cm}
\setlength{\leftmargin}{\evensidemargin}
\addtolength{\leftmargin}{\tmplength}
\settowidth{\labelsep}{X}
\addtolength{\leftmargin}{\labelsep}
\setlength{\labelwidth}{\tmplength}
}
\label{ok_unicode_identifiers_utf8.TfrmArbeitsÜbersicht-pmEntLöchen}
\index{pmEntLöchen}
\item[\textbf{pmEntLöchen}\hfill]
\ifpdf
\begin{flushleft}
\fi
\begin{ttfamily}
public pmEntLöchen: TPopupMenu;\end{ttfamily}

\ifpdf
\end{flushleft}
\fi


\par  \end{list}
\end{document}
